\documentclass[11pt]{article}
\usepackage[a4paper,margin=1in]{geometry}
\usepackage{amsmath,amssymb,amsthm,mathtools}
\usepackage{graphicx}
\usepackage[hidelinks]{hyperref}
\usepackage{microtype}

\newtheorem{theorem}{Theorem}
\newtheorem{lemma}{Lemma}
\newcommand{\HH}{\mathbb H}
\newcommand{\spalg}{\mathfrak{sp}}
\newcommand{\zz}{\mathfrak z}

\title{The exceptional \(C_3\) orbital-measure base pair is singular}
\author{Full answer to an open base case in arXiv:2110.06346}
\date{17 August 2026}

\begin{document}
\maketitle

\paragraph{Status.}
Candidate full negative answer to the explicitly undecided \(C_3\) base
case, with a sharp codimension statement. This is a new-result claim
subject to human review.

\section*{Source question}

Hare's paper asks whether the exceptional pair
\[
       (C_2\times C_1,\; C_1\times C_1)
       \qquad\text{in type }C_3
\]
is absolutely continuous. This is the base case for one of the exceptional
families that become entwined under the proposed rank-reduction procedure.
The question appears on source PDF page~18.

\begin{center}
\includegraphics[width=\textwidth]{figures/open_problem_crop.png}
\end{center}

\section*{Set-up}

Use the compact simply connected model
\[
 G=\operatorname{Sp}(3)
 =\{U\in M_3(\HH):U^*U=I\},\qquad
 \mathfrak g=\spalg(3)=\{A:A^*=-A\}.
\]
Matrices act on the right quaternionic Hilbert space \(\HH^3\).
For \(z\in G\), write \(\zz_z\) for its Lie centralizer and
\(\mathcal N_z=\zz_z^\perp\).

We may take
\[
 x=\operatorname{diag}(1,1,-1),\qquad
 y=\operatorname{diag}(1,-1,e^{i\theta}),\quad 0<\theta<\pi .
\]
Then \(x\) has type \(C_2\times C_1\), while \(y\) has type
\(C_1\times C_1\); the omitted third factor is the empty root system
\(SU(1)\), whose centralizer contributes \(U(1)\).
Thus
\[
\begin{aligned}
\zz_x&\cong\spalg(2)\oplus\spalg(1),\\
\zz_y&=
\{\operatorname{diag}(a_1,a_2,a_3):
 a_1,a_2\in\operatorname{Im}\HH,\ a_3\in\mathbb R i\}.
\end{aligned}
\]

The geometric criterion in source Proposition~2.4 says that
\(\mu_x*\mu_y\) is absolutely continuous exactly when some \(g\in G\)
satisfies
\[
 \mathcal N_x+\operatorname{Ad}(g)\mathcal N_y=\mathfrak g.
\]
Taking orthogonal complements turns this into
\[
 \zz_x\cap\operatorname{Ad}(g)\zz_y=\{0\}.                 \tag{1}
\]

\section*{Proof intuition}

The \(C_2\times C_1\) centralizer is precisely the algebra of
skew-Hermitian operators preserving one quaternionic line \(W\).
After an arbitrary conjugation, the other centralizer is diagonal in three
quaternionic lines; its first two diagonal entries may point along any
imaginary quaternion, while the third is restricted to one imaginary axis.
Given the coordinates of a generator of \(W\), choose a common right
eigenvalue \(q\) so that the third diagonal entry lies on that axis, and
conjugate \(q\) separately by the first two coordinates. This always
manufactures a nonzero diagonal operator that preserves \(W\), so (1) can
never hold.

\section*{Main result}

\begin{theorem}
Let \(x,y\in\operatorname{Sp}(3)\) have types \(C_2\times C_1\) and
\(C_1\times C_1\), respectively. Then
\[
             \mu_x*\mu_y
\]
is purely singular with respect to Haar measure. More sharply, the product
of conjugacy classes \(C_xC_y\) has exact real dimension \(20\) inside the
\(21\)-dimensional group \(\operatorname{Sp}(3)\).
The same conclusion holds in every finite central quotient of
\(\operatorname{Sp}(3)\).
\end{theorem}

\begin{proof}
Conjugate \(x\) to the displayed normal form and put
\(W=e_3\HH\). Fix arbitrary \(g\in G\), let \(v_j=ge_j\), and write a
nonzero generator \(w=e_3\) of \(W\) as
\[
                   w=v_1r_1+v_2r_2+v_3r_3,
                   \qquad r_j\in\HH .
\]
Relative to \(v_1,v_2,v_3\), the algebra
\(\operatorname{Ad}(g)\zz_y\) consists of diagonal operators with
\[
 a_1,a_2\in\operatorname{Im}\HH,\qquad a_3\in\mathbb Ri.  \tag{2}
\]

Suppose first that \(r_3\ne0\). Define
\[
 q=r_3^{-1}ir_3,\qquad a_3=i,\qquad
 a_j=
 \begin{cases}
   r_jqr_j^{-1},&r_j\ne0,\\
   0,&r_j=0
 \end{cases}
 \quad(j=1,2).
\]
Quaternionic conjugation preserves \(\operatorname{Im}\HH\), so (2) holds.
Moreover \(a_jr_j=r_jq\) for every \(j\). Hence the corresponding
diagonal operator \(T\) is nonzero and
\[
                         Tw=wq\in W.                       \tag{3}
\]

If \(r_3=0\), choose any nonzero \(q\in\operatorname{Im}\HH\), put
\(a_3=0\), and use the same formula for \(a_1,a_2\). At least one of
\(r_1,r_2\) is nonzero, so again \(T\ne0\) and (3) holds.

In both cases \(T\in\operatorname{Ad}(g)\zz_y\). Since \(T\) is
skew-Hermitian and preserves \(W\), it also preserves \(W^\perp\):
for \(u\perp W\) and \(w'\in W\),
\[
                   \langle Tu,w'\rangle
                   =-\langle u,Tw'\rangle=0.
\]
Thus \(T\) is block diagonal for \(W^\perp\oplus W\), so
\(T\in\zz_x\). We have proved
\[
 \zz_x\cap\operatorname{Ad}(g)\zz_y\ne\{0\}
 \quad\text{for every }g\in G.
\]
Condition (1) always fails, and the source criterion proves singularity.

It remains to prove the sharp dimension assertion. Take the real
orthogonal, hence quaternionic unitary, matrix
\[
 O=\frac1{9}
 \begin{pmatrix}
  7&-4&-4\\
 -4& 1&-8\\
 -4&-8& 1
 \end{pmatrix}.
\]
All three coordinates \(r_j=O_{3j}\) of \(e_3\) in the column basis of
\(O\) are nonzero real numbers. If
\(T\in\zz_x\cap\operatorname{Ad}(O)\zz_y\), preservation of \(W\) gives
\(Tw=wq\). Comparing the three nonzero real coordinates yields
\(a_1=a_2=a_3=q\); since \(a_3\in\mathbb Ri\), necessarily
\(q\in\mathbb Ri\). Conversely these scalar choices work because \(O\)
is real. Therefore
\[
       \zz_x\cap\operatorname{Ad}(O)\zz_y=\mathbb R(iI),
\]
so the minimum intersection dimension is exactly one.
Since \(\dim_{\mathbb R}\spalg(3)=21\), the maximum differential rank of
the conjugacy-class product map is \(21-1=20\). It is attained at \(O\);
the constant-rank theorem locally and the universal upper bound globally
give \(\dim(C_xC_y)=20\). Finite central quotients preserve this local
dimension and the interior/absolute-continuity criterion.
\end{proof}

\section*{Verification}

The included script \texttt{code/verify\_c3\_pair.py} performs two independent
checks. First, it tests the quaternionic construction in \(20{,}000\)
random coefficient triples, including purity of the \(a_j\) and every
identity \(a_jr_j=r_jq\). Second, it computes over the rationals that the
six off-block constraints induced by \(O\) have rank six on the seven
parameters of \(\zz_y\), with kernel
\[
                         (1,0,0,1,0,0,1),
\]
confirming the one-dimensional intersection and rank \(20\).

\section*{Upgrade attempts and limitations}

The first upgrade sharpened singularity to exact codimension one.
A second focused attempt tried to propagate the same construction to the
higher exceptional \(C_n\) families. It does not do so: in higher rank the
smaller \(x\)-eigenspace is an \(m\)-plane rather than a quaternionic line,
and preserving that plane imposes a matrix invariance condition not solved
by the single right-eigenvalue construction. Dimension count permits zero
intersection there. Accordingly, this packet does not claim the general
\(C_n,D_n\) conjecture or the higher exceptional families.

\section*{Novelty check and review recommendation}

A bounded search covered the run indexes, arXiv:2110.06346, exact searches
for the displayed \(C_3\) type pair and its absolute continuity, later
orbital-measure work by the source author, and searches combining the terms
exceptional, orbital measure, and \(C_3\). No later resolution of this base
case was found. Novelty confidence is moderate, not exhaustive.

Human review should focus on the identification of the two quaternionic Lie
centralizers and on the orthogonal-complement form of source
Proposition~2.4. Once those standard identifications are accepted, the
intersection construction and the sharpness computation are elementary.

\section*{Reference}

\begin{sloppypar}
K. E. Hare, \emph{The absolute continuity of convolutions of orbital
measures in \(SO(2n+1)\)}, arXiv:2110.06346 (2021), especially
Proposition~2.4 and the discussion preceding Conjecture~5.9.
\end{sloppypar}

\end{document}
